% -----------------------------------------------------------------------------
% Comandos científicos
% -----------------------------------------------------------------------------
\newcommand{\channelset}{\mathcal{S}}
\newcommand{\subsetSI}{S_I}
\newcommand{\loss}{\mathcal{L}}
\newcommand{\meanloss}[3][]{\overline{L}^{#1}_{#2,#3}(n)}
\newcommand{\Nstar}{\ensuremath{N^{\star}}}
\newcommand{\rankrho}{\ensuremath{\rho_{\mathrm{rank}}}}
\newcommand{\WinnersAgreement}{\textit{Winners Agreement}}
\newcommand{\NstarSimilarity}{\textit{Progressive }\ensuremath{N^{\star}}\textit{-Similarity}}
\newcommand{\RealToReal}{\ensuremath{\text{Real}\rightarrow\text{Real}}}
\newcommand{\GaussianToReal}{\ensuremath{\text{Gaussiana}\rightarrow\text{Real}}}
\newcommand{\GMMToReal}{\ensuremath{\text{GMM}\rightarrow\text{Real}}}
\newcommand{\CoInfoSim}{\textsf{CoInfoSim}}
\newcommand{\SLACGS}{\textsf{SLACGS}}

% -----------------------------------------------------------------------------
% Fonte de ilustrações, tabelas e quadros: obrigatória inclusive para autoria própria.
% -----------------------------------------------------------------------------
\newcommand{\sourcebelow}[1]{%
  \par\vspace{2pt}%
  {\sffamily\footnotesize Fonte: #1.\par}
}
\newcommand{\notailustracao}[1]{%
  \par{\sffamily\footnotesize Nota: #1.\par}
}

% -----------------------------------------------------------------------------
% Caixas editoriais reutilizáveis. Estilo contido: fundo claro, filete de
% destaque à esquerda, sem molduras pesadas. Todas legíveis em tons de cinza.
% -----------------------------------------------------------------------------
\tcbset{coinbox/.style={
  enhanced,
  breakable,
  boxrule=0pt,
  frame hidden,
  arc=1mm,
  left=3mm,right=3mm,top=2mm,bottom=2mm,
  borderline west={2.2pt}{0pt}{CoinAccent},
  fonttitle=\bfseries\sffamily\small,
  coltitle=CoinAccentDark,
  attach title to upper={\par\smallskip},
}}

% Definição formal ou operacional de um conceito.
\newtcolorbox{definitionbox}[1][]{
  coinbox,
  colback=CoinAccentPale,
  title={#1},
}

% Interpretação: o que um resultado significa e o que não demonstra.
\newtcolorbox{interpretationbox}[1][Interpretação]{
  coinbox,
  colback=CoinGrayPale,
  borderline west={2.2pt}{0pt}{CoinGray},
  coltitle=black,
  title={#1},
}

% Síntese de resultado ao final de subseções de resultados.
\newtcolorbox{keyresultbox}[1][Síntese do resultado]{
  coinbox,
  colback=CoinAccentPale,
  title={#1},
}

% Registro destacado de limitação ou ameaça à validade.
\newtcolorbox{limitationbox}[1][Limitação]{
  coinbox,
  colback=CoinGrayPale,
  borderline west={2.2pt}{0pt}{CoinGray},
  coltitle=black,
  title={#1},
}

% Pergunta de pesquisa.
\newtcolorbox{researchquestionbox}[1][]{
  coinbox,
  colback=white,
  colframe=white,
  title={#1},
}

% -----------------------------------------------------------------------------
% Marcadores de conteúdo pendente. Somente aparecem no modo rascunho.
% -----------------------------------------------------------------------------
\newcommand{\editorialtodo}[1]{%
  \ifdraftdocument
    \begin{editorialbox}#1\end{editorialbox}
  \fi
}

\newcommand{\figureplaceholder}[3][0.82\textwidth]{%
  \begin{figure}[H]
    \centering
    \caption{#2}
    \fbox{\parbox[c][0.24\textheight][c]{#1}{\centering\sffamily\small #3}}
    \sourcebelow{elaboração própria; figura definitiva a inserir a partir dos relatórios publicados do \CoInfoSim}
  \end{figure}
}

\newcommand{\matrixplaceholder}[3][0.46\textwidth]{%
  \begin{minipage}[t]{#1}
    \centering
    \fbox{\parbox[c][0.20\textheight][c]{0.94\linewidth}{\centering\sffamily\small #3}}
    \subcaption{#2}
  \end{minipage}
}
